% META: {"id":"1SPE-GEOREP-CO-012","chapitre":"1SPE-GEOMETRIE-REPEREE","type_objet":"corrige","exercice_ref":"1SPE-GEOREP-EX-012","capacites_codes":["C2"],"sources_inspiration":[],"status":"approved","origin":"generated","status_history":["generated","audited","accepted","approved"],"release_acceptance":"RELEASE_OWNER_BATCH_ACCEPTANCE","acceptance_closure_digest":"sha256:9b3ccf9a81c5520fbb7e03b7b2d3e7bf2057a02834b6d908bf3be5f49f3b553f","acceptance_receipt_digest":"sha256:4fad86f0282e0fa4e0419cca1ad6379ae7af5a280c04a4a90880f90a1c044242"}
% BEGIN-VERIFY
% from sympy import *
% x, y = symbols('x y')
% # Droite passant par A(3,-2) de vecteur normal n(1,4)
% # 1(x-3) + 4(y+2) = 0 => x + 4y + 5 = 0
% assert 1*3 + 4*(-2) + 5 == 0
% END-VERIFY
\begin{corrige}{1SPE-GEOREP-EX-012}

\commentaireMarge{On utilise la formule $a(x - x_0) + b(y - y_0) = 0$.}

$1(x - 3) + 4(y + 2) = 0$, soit $x + 4y + 5 = 0$.

Vérification : $3 + 4 \times (-2) + 5 = 3 - 8 + 5 = 0$ \checkmark.

\end{corrige}
